\textbf{Proof.}
Let \(\mathcal C\) be the sigma-field generated by all source and target covariates. The corrected spline contract changes only the finite dimensions: the interior regression basis has \(d_{\mu,n}=J_{\mu,n}+q\) coordinates, and the nonconstant calibration vector has \(d_{r,n}-1=J_{r,n}+q-1\) coordinates. The positive ridge terms in \(\mathrm{def:regression\text{-}fit}\) make every foldwise regression coefficient a unique linear function of the source pseudo-outcomes in its training complement. The corrected matrices and smoother coefficients in \(\mathrm{def:ricot\text{-}linear\text{-}coefficients}\) are exactly the coefficients obtained from the corresponding normal equations. Therefore, for every fold \(k\) and every \(x\in[0,1]\),
\[
\widehat\mu^{(-k)}_{h,J_{\mu,n}}(x)
=\sum_{q_0=1}^n L_{q_0,k,h,J_{\mu,n}}(x)Z_{q_0}.
\tag{1}
\]
Substitution of (1) into the target plug-in term and the fitted-regression term of \(\widehat\tau^{\mathrm{sp}}_{h,J_{\mu,n},J_{r,n}}\) shows that the coefficient of a fixed \(Z_{q_0}\) is
\[
\begin{aligned}
{}&\frac1m\sum_{j=1}^mL_{q_0,k_T(j),h,J_{\mu,n}}(X_j^T)
+\frac1n\widehat r^{\mathrm{RICOT},(-k_S(q_0))}_{h,J_{r,n}}(X_{q_0})\mathbf 1\{X_{q_0}\ge h\}\\
&\quad-\frac1n\sum_{i=1}^n
\widehat r^{\mathrm{RICOT},(-k_S(i))}_{h,J_{r,n}}(X_i)\mathbf 1\{X_i\ge h\}
L_{q_0,k_S(i),h,J_{\mu,n}}(X_i).
\end{aligned}
\tag{2}
\]
Expression (2) is \(\Gamma_{q_0}(h,J_{\mu,n},J_{r,n})\) by \(\mathrm{def:ricot\text{-}linear\text{-}coefficients}\), proving the first identity. Every factor in (2) is \(\mathcal C\)-measurable: the ridge Gram matrices and the branch-qualified RiCOT primal--dual programs depend only on the observed covariates. The zero-interior-mass fallback is deterministic. On a positive-mass branch the corrected standard and reparameterized spline vectors have the displayed finite dimensions, so their selected optimizers and tie-broken multipliers are measurable finite-dimensional functions of the covariates. The pointwise cap is measurable as well. Hence every \(\Gamma_{q_0}\) is \(\mathcal C\)-measurable.

By consistency, conditional trial randomization, and the definition of the known randomization propensity, for \(P_S\)-almost every \(x\),
\[
\begin{aligned}
\mathbb E_{P_S}(Z\mid X=x)
&=\frac{\mathbb E_{P_S}(AY\mid X=x)}{e(x)}
-\frac{\mathbb E_{P_S}\{(1-A)Y\mid X=x\}}{1-e(x)}\\
&=\mathbb E_{F^S_{n,m}}\{Y^S(1)\mid X=x\}
-\mathbb E_{F^S_{n,m}}\{Y^S(0)\mid X=x\}\\
&=\mu_1(x)-\mu_0(x)=\mu(x).
\end{aligned}
\tag{3}
\]
The divisions in (3) are valid because treatment positivity gives \(e(x)\in[\varepsilon,1-\varepsilon]\). Adding and subtracting \(\sum_i\Gamma_i\mu(X_i)\) and \(m^{-1}\sum_j\mu(X_j^T)\) in the first identity gives the displayed exact error decomposition.

The H\"older norm is subadditive. Thus the two response-smoothness assumptions imply \(\mu=\mu_1-\mu_0\in\mathcal H^s(2M)\). Since \(0\le\mu_d\le1\), one also has \(\|\mu\|_\infty\le1\). The definition of \(B_X\) consequently yields the designwise bound
\[
\left|\sum_{i=1}^n\Gamma_i\mu(X_i)-\frac1m\sum_{j=1}^m\mu(X_j^T)\right|\le B_X.
\tag{4}
\]

Treatment positivity and \(0\le Y\le1\) imply \(-\varepsilon^{-1}\le Z_i\le\varepsilon^{-1}\). Conditional on \(\mathcal C\), the source pseudo-outcomes are independent and centered as in (3). The conditional range length of \(\Gamma_i\{Z_i-\mu(X_i)\}\) is at most \(2|\Gamma_i|/\varepsilon\). Applying \(\mathrm{lem:bounded\text{-}sum\text{-}hoeffding\) conditionally with \(\delta=\gamma/2\) gives
\[
\mathbb P_{n,m}\!\left(
\left|\sum_{i=1}^n\Gamma_i\{Z_i-\mu(X_i)\}\right|>R_S
\,\middle|\,\mathcal C\right)\le\frac\gamma2.
\tag{5}
\]
Independence of the target sample, the identity \(\mathbb E_{P_T}\mu(X^T)=\tau\), and the bound \(\mu(X^T)\in[-1,1]\) permit a second application of the same lemma to the independent variables \(m^{-1}\{\mu(X_j^T)-\tau\}\). Their range lengths are \(2/m\). Again using \(\delta=\gamma/2\),
\[
\mathbb P_{n,m}\!\left(
\left|\frac1m\sum_{j=1}^m\mu(X_j^T)-\tau\right|>R_T(m)
\right)\le\frac\gamma2.
\tag{6}
\]
Taking expectations in (5), combining (4)--(6) with the exact decomposition, and applying the union bound prove the coverage probability. Finally, \(B_X\le b\) and \(\sum_i\Gamma_i^2\le g^2\), together with the definitions of \(R_S\) and \(R_T\), give the asserted realized-design radius. No density-smoothness condition, density-ratio localization, leverage bound, or stochastic rate for the two observable diagnostics enters the argument. \(\square\)